\documentclass[a4paper,11pt]{ltjsarticle}
\newcommand{\bookletname}{問題冊子}
% LuaLaTeX / UTF-8
\usepackage{amsmath,amssymb,bm}
\usepackage[margin=23mm,top=22mm,bottom=22mm]{geometry}
\usepackage{enumitem,booktabs,array}
\usepackage{fancyhdr}
\usepackage[hidelinks,unicode]{hyperref}
\setlength{\parindent}{1\zw}
\setlength{\parskip}{0.35em}
\setlength{\headheight}{17pt}
\pagestyle{fancy}
\fancyhf{}
\fancyhead[L]{\small 第1回模試\quad 統計応用（医薬生物学）}
\fancyhead[R]{\small\bookletname}
\fancyfoot[C]{\small\thepage}
\renewcommand{\headrulewidth}{0.35pt}
\setlist[enumerate,1]{label=〔\arabic*〕,leftmargin=3.5em,labelsep=0.5em,itemsep=0.85em,topsep=0.7em}
\newcommand{\E}{\operatorname{E}}
\newcommand{\Var}{\operatorname{Var}}
\newcommand{\Cov}{\operatorname{Cov}}
\newcommand{\Prob}{\operatorname{P}}
\newcommand{\qhead}[1]{\section*{問#1}}
\newcommand{\ans}[1]{\subsection*{〔#1〕}}
\newcommand{\points}[1]{\hfill\textnormal{（#1点）}}
\newcommand{\continued}[1]{\noindent\textbf{問#1（続き）}\par\medskip}
\hypersetup{pdfauthor={Codex},pdfsubject={統計検定1級対策・自主学習用模擬試験}}

\hypersetup{pdftitle={第1回模試 統計応用（医薬生物学） 問題冊子}}
\begin{document}
\thispagestyle{empty}
\begin{center}
\vspace*{10mm}
{\large 2026年 統計検定1級対策}\par\vspace{7mm}
{\Huge\bfseries 第1回 模擬試験}\par\vspace{6mm}
{\LARGE 統計応用（医薬生物学）}\par\vspace{6mm}
{\Large 問題冊子}\par\vspace{10mm}
{\large 試験時間 90分\qquad 5問中3問選択}
\end{center}
\vspace{6mm}
\noindent\textbf{受験上の注意}
\begin{enumerate}[label=\arabic*.,leftmargin=2em,itemsep=0.65em]
\item 問1から問5のうち3問を選び，選択した問題番号を答案に明記すること。解答は別紙に記述すること。
\item 答えのみでよいと明記した箇所を除き，導出または理由を示すこと。各問の記号はその問題の中だけで用いる。
\item 四則演算・平方根が計算できる一般電卓と，別添の「統計検定1級\_付表数値表」を使用してよい。付表1～5の全体を事前に用意すること。本冊子には付表を再録していない。
\item $\log$ は自然対数を表す。数値は特に指定がなければ小数第3位程度まで求め，付表の丸めや補間による差は許容する。
\item 本模試では各大問を30点とする。選択した3問の合計90点を100点に換算する。小問の配点は採点基準冊子に記載する。
\item 問5は統計応用の共通問題を想定した問題である。データ・研究設定はすべて本模試のための架空のものである。
\end{enumerate}
\vfill
\begin{center}
氏名\ \underline{\hspace{55mm}}\qquad 実施日\ \underline{\hspace{35mm}}\par\vspace{8mm}
\small 自主学習用のオリジナル模試です。公式の問題・配点ではありません。\par
作成日：2026年11月12日\qquad 作成識別：6astra
\end{center}
\newpage

\qhead{1}
ある疾患の発症から死亡までの時間を比較する観察研究で，対照群を $x=0$，治療群を $x=1$ とする。
発症時に登録される人と，発症後しばらく経ってから登録される人がいる。
発症を時間0とした登録時点 $a_i$，観察終了時点 $y_i$，死亡指標 $\delta_i$ を表に示す。
$\delta_i=0$ は右側打ち切りを表し，登録前に死亡した人は観察されない。
個人間は独立であり，発症者の母集団では各群内で登録時点および打ち切りの仕組みは死亡時間と独立と仮定する。登録時点・打ち切り時点の分布は両群で共通とする。
死亡時点 $t$ のリスク集合には $a_i\leq t\leq y_i$ を満たす人だけを含める。
\begin{center}
\begin{tabular}{ccrrc}
\toprule
被験者 & $x_i$ & $a_i$（月） & $y_i$（月） & $\delta_i$\\
\midrule
C1 & 0 & 0 & 1 & 1\\
C2 & 0 & 0 & 3 & 0\\
C3 & 0 & 2.5 & 5 & 1\\
C4 & 0 & 0 & 6 & 0\\
T1 & 1 & 0 & 2 & 1\\
T2 & 1 & 0 & 4 & 1\\
T3 & 1 & 3.5 & 6 & 0\\
T4 & 1 & 0 & 6 & 0\\
\bottomrule
\end{tabular}
\end{center}
\begin{enumerate}
\item 各死亡時点について，対照群と治療群のリスク集合の人数を示せ。カプラン・マイヤー法による，各群の $t=5$ における生存確率の推定値を求めよ。
\item Cox比例ハザードモデル
\[
\lambda(t\mid x)=\lambda_0(t)e^{\beta x},\qquad \beta\in\mathbb R
\]
を当てはめる。$\theta=e^\beta$ とおき，部分尤度を求め，これを最大にするハザード比 $\hat\theta$ を求めよ。最大であることも確認すること。
\item 部分対数尤度を $\ell(\beta)$ とする。$U(0)=\ell'(0)$ と $I(0)=-\ell''(0)$ を求めよ。
帰無仮説 $H_0:\beta=0$ の下で $U(0)^2/I(0)$ が近似的に自由度1のカイ二乗分布に従うとして，両側5\%で検定せよ。$U(0)$ とログランク検定のスコアとの関係も述べよ。
\item $a_i$ を無視し，全員を時間0から観察したものとして，両群を併合した生存関数を推定する担当者がいた。この方法と正しいリスク集合を使う方法で，それぞれ $t=5$ のカプラン・マイヤー推定値を計算せよ。差が生じる理由を述べよ。
\end{enumerate}
\newpage

\qhead{2}
2地域A，Bの死亡率を比較する。年齢を若年層と高年層に分け，観察人年と死亡数を表に示す。
基準集団の年齢別死亡率は既知であり，観察人年は固定値として扱う。
\begin{center}
\begin{tabular}{lrrrrr}
\toprule
 & 基準死亡率 & \multicolumn{2}{c}{地域A} & \multicolumn{2}{c}{地域B}\\
\cmidrule(lr){3-4}\cmidrule(lr){5-6}
年齢層 & （1人年当たり） & 観察人年 & 死亡数 & 観察人年 & 死亡数\\
\midrule
若年 & 0.010 & 5000 & 40 & 1250 & 10\\
高年 & 0.040 & 1250 & 60 & 5000 & 240\\
\bottomrule
\end{tabular}
\end{center}
基準死亡率を地域の各年齢層の観察人年に適用した期待死亡数を $E_k$ とする。
標準化死亡比（SMR）は，死亡総数を $\sum_kE_k$ で割ったものと定義する。
\begin{enumerate}
\item 地域A，Bの粗死亡率とSMRを求めよ。
\item 両年齢層の重みをそれぞれ $1/2$ とする直接標準化死亡率を，地域A，Bについて求めよ。
一般に，地域の年齢別死亡率が基準死亡率 $r_k$ の $\theta_k$ 倍で，観察人年が $Y_k$ のとき，SMRの母数に対応する量
\[
\frac{\sum_kY_kr_k\theta_k}{\sum_kY_kr_k}
\]
を $\theta_k$ の重み付き平均として表せ。これを用いて，この表で両地域のSMRが異なる理由を述べよ。
\item 以下では地域Aだけを考える。独立な死亡数 $O_k$ が平均 $\theta E_k$ のポアソン分布に従うモデルを仮定する（$\theta\geq0$，漸近近似は $\theta>0$ で考える）。$\theta$ の最尤推定量，その期待値と分散を求めよ。
さらにデルタ法で $\log\hat\theta$ の漸近分散を求め，対数尺度の正規近似を用いた $\theta$ の95\%信頼区間を数値で求めよ。
\item 〔3〕の共通の $\theta$ を仮定するモデルに対し，Pearson型適合度統計量
\[
X^2=\sum_{k=1}^{2}\frac{(O_k-\hat\theta E_k)^2}{\hat\theta E_k}
\]
を計算せよ。その漸近帰無分布の自由度を理由とともに答え，有意水準5\%で適合度を検定せよ。
\item 〔1〕～〔4〕を踏まえ，「地域AのSMRは1なので，どの年齢層でも基準集団と死亡率が等しい」「SMRが大きい地域Bでは，同じ年齢の住民でも地域Aより死亡率が高い」という2つの主張を評価せよ。
\end{enumerate}
\newpage

\qhead{3}
ある疾患のスクリーニング検査を評価する。検査結果を $S$（陽性なら1，陰性なら0），真の疾患の有無を $D$（疾患ありなら1）とする。
対象母集団における陽性割合 $q=\Prob(S=1)=0.20$ は既知とする。
陽性者200人，陰性者800人からなる評価用集団において，陽性者は全員，陰性者は無作為に選んだ100人だけに確定診断を行った。その結果を表に示す。確定診断には誤りがないものとする。
\begin{center}
\begin{tabular}{lrrr}
\toprule
 & 確定診断した人数 & 疾患あり & 疾患なし\\
\midrule
検査陽性 & 200 & 80 & 120\\
検査陰性 & 100 & 5 & 95\\
\bottomrule
\end{tabular}
\end{center}
$a=\Prob(D=1\mid S=1)$，$b=\Prob(D=1\mid S=0)$ とおく。
確定診断された疾患ありの人数 $X_1,X_0$ は，それぞれ独立に
$\operatorname{Bin}(200,a)$，$\operatorname{Bin}(100,b)$ に従うとする。
\begin{enumerate}
\item 確定診断された300人だけをそのまま集計して求めた感度と特異度の値を答えよ。この集計で検査陽性者の構成割合はどのようになっているか。
\item 陰性者の抽出割合を考慮して，評価用集団1000人の $S\times D$ 表の各セルの人数を推定せよ。この表から，感度・特異度・陽性的中率・陰性的中率を推定せよ。
\item 母集団の有病割合 $p=\Prob(D=1)$ を $q,a,b$ で表せ。$q$ を既知とする不偏推定量 $\hat p$ を構成し，その分散を求めよ。未知母数を標本割合で置き換えた標準誤差と，正規近似による $p$ の95\%信頼区間を数値で求めよ。
\item 一般に，確定診断の実施指標を $R$ とし，$S=s$ の下で $R$ と $D$ が独立で，$\Prob(R=1\mid S=s)=v_s>0$ とする。
真の感度を $h=\Prob(S=1\mid D=1)$ とするとき，確定診断された人の中での感度 $\Prob(S=1\mid D=1,R=1)$ を $h,v_1,v_0$ で表せ。
さらに，陽性的中率はこの抽出により変わるかを示せ。陰性者の中でも疾患が疑われる人ほど確定診断を受けやすい場合，〔2〕の補正をそのまま使えるか，理由を述べよ。
\end{enumerate}
\newpage

\qhead{4}
施設を単位として新しい介入を無作為に割り付ける試験を考える。対照群を $g=0$，介入群を $g=1$ とし，各群 $K$ 施設，各施設 $m$ 人を観察する。改善量 $Y_{gij}$ は
\[
Y_{gij}=\mu+\tau g+B_{gi}+\varepsilon_{gij},\qquad
i=1,\ldots,K,\quad j=1,\ldots,m
\]
に従うとする。ここで $B_{gi}\sim N(0,\sigma_B^2)$，$\varepsilon_{gij}\sim N(0,\sigma_e^2)$ はすべて互いに独立で，$\sigma_B^2,\sigma_e^2>0$ は既知とする。
$\tau$ が大きいほど介入の効果が大きい。
\begin{enumerate}
\item 個人の測定値の分散，同じ施設の異なる2人の測定値の共分散と相関係数 $\rho$ を求めよ。
\item 各群の全員の平均を $\bar Y_g$ とする。$\hat\tau=\bar Y_1-\bar Y_0$ の期待値と分散を求めよ。
全員を独立と誤認し，各個人の周辺分散を用いて計算した $\hat\tau$ の分散を $V_{\mathrm{ind}}$ とするとき，$\Var(\hat\tau)/V_{\mathrm{ind}}$ を $m,\rho$ で表せ。
\item $K=12$，$m=10$，$\sigma_B^2=4$，$\sigma_e^2=36$，観測された差が $\hat\tau=3$ であった。
$H_0:\tau=0$ を $H_1:\tau\ne0$ に対して有意水準5\%で検定し，$\tau$ の95\%信頼区間を求めよ。
\item 同じ分散成分を用いて，$m=10$ の試験を新たに計画する。
$H_0:\tau\leq0$ を $H_1:\tau>0$ に対して片側2.5\%で検定し，$\tau=2$ のとき検出力80\%以上となる最小の各群施設数 $K$ を求めよ。
標準正規分布の上側20\%点には付表から得る $0.84$ を用いてよい。
\item 一施設を参加させる費用を $c_0>0$，一人を追加観察する費用を $c_1>0$ とする。
各群の予算を $C$ とし，$K(c_0+c_1m)=C$ を満たすように計画する。
まず $K,m$ を正の実数として，$\Var(\hat\tau)$ を最小にする $m$ を導出せよ。
$\sigma_B^2=4,\sigma_e^2=36,c_0=100,c_1=1,C=5200$ の場合について，最適な $m,K$ を求めよ。
\end{enumerate}
\newpage

\qhead{5（共通問題形式）}
複数の改善案を別々のデータで評価する。案 $j$ の効果の推定値 $X_j$ は独立に
$N(\mu_j,\sigma^2)$ に従い，$\sigma>0$ は既知とする。最も大きい推定値を示した案を選ぶことを考える。
\begin{enumerate}
\item 案が2つで $\mu_1=\mu_2=0$ の場合，$M=\max(X_1,X_2)$ の期待値と分散を求めよ。
答えは $\sigma,\pi$ を含む式でよい。
\item 案が4つで $\mu_1=\cdots=\mu_4=0$ の場合，$M=\max_{1\leq j\leq4}X_j$ の累積分布関数を標準正規分布の累積分布関数 $\Phi$ で表せ。
各案について片側有意水準5\%で $H_j:\mu_j\leq0$ を検定し，1つでも棄却されたら効果のある案が見つかったとする。このときの第1種の過誤確率を求めよ。
\item 4つの仮説全体の家族内第1種過誤率（FWER）を5\%以下にするため，Bonferroni法を用いる。
各 $X_j/\sigma$ に対する共通の棄却限界値を付表から求めよ。
\item 4案の片側検定統計量 $X_j/\sigma$ が，案1から順に $2.58,2.10,1.96,0.20$ であった。
各片側 $p$ 値を求め，Bonferroni法と次のHolm法（ともにFWER 5\%）で棄却される仮説をそれぞれ示せ。

Holm法では $p$ 値を小さい順に $p_{(1)}\leq\cdots\leq p_{(4)}$ と並べる。
$k=1$ から順に $p_{(k)}\leq0.05/(5-k)$ ならその仮説を棄却して次へ進み，初めてこの条件が満たされなければ，そこで止めて残りをすべて棄却しない。
\item 一般に $m$ 個の仮説に対し，Holm法の第 $k$ 段階の基準を $\alpha/(m-k+1)$ とする。
真の帰無仮説の集合を $\mathcal T$，その個数を $m_0\geq1$ とし，それぞれの $p$ 値が
\[
\Prob(p_j\leq u)\leq u\quad(0\leq u\leq1,\ j\in\mathcal T)
\]
を満たすものとする。$p$ 値相互の独立性を仮定せずに，Holm法のFWERが $\alpha$ 以下であることを示せ。
\end{enumerate}
\vfill
\begin{center}（問題は以上）\end{center}
\end{document}
